Este trabajo presenta el desarrollo de un sistema de control semafórico adaptativo basado en \textbf{Deep Reinforcement Learning (DRL)} utilizando el algoritmo \textbf{Proximal Policy Optimization (PPO)} integrado con el simulador microscópico de tráfico \textbf{SUMO}. El objetivo principal es reducir la congestión urbana mediante la optimización dinámica de tiempos semafóricos en una intersección crítica.

El agente aprende a manejar condiciones de tráfico altamente variables gracias a una función de recompensa diseñada para balancear eficiencia y estabilidad, incorporando penalizaciones por colas, desequilibrio direccional y cambios excesivos de fase. El entrenamiento se realizó bajo un perfil de tráfico dinámico de 60 minutos con múltiples picos de demanda.

Los resultados muestran una \textbf{reducción aproximada del 60\% en la longitud promedio de colas} y un \textbf{incremento del 4.5\% en el flujo vehicular efectivo (throughput)} respecto a un controlador de tiempo fijo. Estos resultados demuestran la efectividad del enfoque DRL + SUMO para mejorar la movilidad urbana. Finalmente, se discuten las limitaciones actuales y se proponen líneas de trabajo futuro, incluyendo control multiagente y percepción mediante visión por computador.

\vspace{1em}
\noindent \textbf{Palabras clave:} semáforos inteligentes, aprendizaje por refuerzo, PPO, SUMO, tráfico urbano.